\chapter{绪论}
本章介绍注视点渲染技术的研究背景与意义，阐述当前国内外研究现状，并简要说明本文的研究内容和结构安排。

\section{研究背景与意义}

虚拟现实（Virtual Reality VR）技术近年来迅速发展，通过模拟真实世界的视觉、听觉、触觉等多感官体验，使用户沉浸于计算机生成的虚拟三维环境中。
这种沉浸式体验使用户能够更加深度地参与虚拟环境，并与之进行互动。当前，虚拟现实技术已广泛应用于多个领域，如游戏娱乐、教育、医疗、军事等。
在游戏娱乐领域，虚拟现实技术为玩家提供了更加真实和富有沉浸感的游戏体验，推动了游戏产业的革新与发展。
在教育领域，虚拟现实能够模拟各种实际中难以实现的环境，帮助学生更加直观地理解和掌握复杂的知识与技能。
在医疗领域，虚拟现实被广泛应用于手术模拟训练、康复治疗等方面，提高了医疗操作的安全性和效率[3]。
在军事领域，虚拟现实技术被用来构建逼真的战斗模拟场景，提升士兵的战斗技能及应急反应能力[4]。
此外，虚拟现实能够将现实世界中高昂或难以承担的风险转化为虚拟风险，从而为决策者提供宝贵的经验与洞察，应用于现实世界中的训练和决策过程。

在大规模虚拟现实场景中，尤其涉及复杂动画与角色的渲染时，用户对虚拟场景的真实性和交互实时性要求不断提高，这就意味着动画网格的复杂度逐渐增加，
给人以真实感的同时带来了更大的数据量，其中包括大量的顶点数据、纹理数据、骨骼数据及动画序列。
实现这种高真实性与高实时性的渲染通常需要消耗大量的计算资源和时间，从而显著增加了渲染延迟和存储开销。
因此，虚拟现实中动画网格的渲染着面临严峻挑战：

（1）动画网格数据规模过大显著增加渲染延迟和存储开销[5]。在虚拟现实的渲染流水线中，CPU和GPU的通信是至关重要的一步，具体流程如下：首先，CPU负责构建动画网格，包括网格数据（顶点、法线、纹理坐标等）和动画数据（如骨骼和关键帧）。然后，CPU对骨骼变换和顶点蒙皮进行计算，优化数据并将其传输到GPU。GPU端利用顶点着色器处理顶点变换和蒙皮，应用骨骼矩阵，执行几何着色器，进行光栅化，计算光照和纹理效果，最终CPU调用渲染命令让GPU执行渲染任务，通过片段着色器生成像素颜色，并输出到帧缓冲区，从而生成最终的二维图像。然而，渲染命令数量的增加会导致性能瓶颈，尤其是在处理复杂且精细的动画网格时，这些网格通常包含庞大的网格数据和动画数据，渲染过程中需要频繁调用渲染命令，从而加重了渲染延迟和存储开销。目前，传统的多细节层次优化（LOD）方法通常根据视距动态调整网格的细节层级，旨在保持视觉效果一致的情况下减少网格细节，但该方法仅限于静态网格细节的简化，未能有效考虑动画序列的优化。因此，针对虚拟现实中的复杂动画网格渲染问题，需要设计一种多维层次模型，旨在精简静态网格数据并优化动态动画序列。该模型应能够适应不同分辨率的渲染要求，同时有效降低渲染延迟和存储开销。

（2）注视点渲染忽略视觉锐度的时间性变化导致渲染延迟。人类视觉系统 （HVS） 能够在中央凹（人类视网膜中心）区域感知场景中的最精细细节。视觉锐度是眼睛能够辨别细节的能力，通常以光栅分辨率（以周期/度为单位）来衡量。尽管HVS 具有较宽的视野（FOV），但其视觉锐度最高的区域（中央凹区域）仅覆盖视野中央约 2.5°的范围。注视点渲染技术利用这一特性，有选择性地对图形帧中的中心区域进行高精度渲染。现有的注视点渲染算法充分利用了空间感知的变化，但尚未有效利用眼跳（短时间内迅速改变注视焦点）过程的扫视抑制现象及眼跳落点后的短时间内视觉锐度恢复的时间延迟特性[6]。在此期间，中央凹区域的视觉清晰度尚未恢复至最佳水平，因此无需进行高质量渲染。忽视这一特性会导致渲染资源的浪费，从而引发渲染延迟及帧间闪烁现象，严重影响用户的视觉体验和沉浸感。因此，亟需设计一种基于眼跳效应的注视点渲染方法，旨在保持视觉效果的精细度的同时减少不必要的渲染计算，从而优化渲染效率并降低延迟。

本文聚焦于上述挑战，进行了如下研究：
（1）针对动画网格数据规模过大显著增加渲染延迟和存储开销的问题，分析动画网格的组成结构，通过解耦静态网格与动画数据，采用网格合并和纹理压缩技术精简静态网格数据、优化关键帧和精简骨骼数据来优化动画序列，构建适应不同分辨率渲染的多维层次模型，以降低渲染负担，提高渲染效率。

（2）针对注视点渲染中忽略视觉锐度的时间性变化引发的渲染延迟问题，深入研究扫视抑制现象及眼跳后视觉锐度恢复延迟的特性，提出基于眼跳效应的注视点渲染方法，在扫视与眼跳后视觉锐度恢复期间动态降低渲染分辨率，减少潜在计算和数据开销，从而提高渲染效率。

（3）结合上述技术，设计并实现了眼跳引导的虚拟现实中动画网格的高效渲染原型系统，有效缓解了虚拟现实中动画网格的高效渲染的问题，提升了虚拟现实系统的用户体验、虚拟场景的构建真实性及交互实时性。
\section{国内外研究现状}

在上述研究背景下，本节将对动画网格在虚拟现实中的3D场景构建与渲染加速方法进行综述，重点探讨相关技术的研究进展。随着虚拟现实应用的不断发展，如何提升渲染效率、降低计算资源消耗、优化视觉效果成为当前的研究重点。网格简化、贴图压缩和注视点渲染等技术被广泛应用于虚拟现实场景中，旨在平衡渲染质量与性能需求。

\subsection{网格简化研究现状}

网格简化是一种重要的三维优化技术，它通过减少网格的顶点数和三角面片数来降低网格渲染的计算复杂度，以提高渲染效率。在虚拟现实、增强现实、游戏、建筑等领域中，网格简化技术都被广泛应用，为这些领域带来了更高的渲染性能。针对传统的网格简化算法存在的缺陷和不足，研究者提出了许多优化网格简化算法的思路，故如何优化网格简化算法是如今网格简化技术的研究重点。网格简化算法是实现网格简化的关键，如今成熟的网格简化算法有基于拓扑结构优化的网格简化算法、基于误差度量的网格简化算法[7]、基于边坍缩的网格简化算法、基于曲面拟合的网格简化算法[8]、基于自适应的网格简化算法等，研究者也正在探索各种新的算法来提升网格简化的效果。网格简化后的网格质量是影响其渲染效果的重要因素，进行网格简化后还应该对简化后的网格进行质量评估，具体的评估方法有误差分析、形状相似度分析、光照一致性分析等。

作为进行网格简化处理的基础,网格简化算法在实际工程中运用时有一定的开发成本，为了使开发者更简易地进行网格简化处理，如今市面上有许多成熟的网格简化工具可供使用。MeshLab[9]是一款开源免费的网格处理软件，支持各种网格简化算法,并且还提供了可视化编辑界面,能够方便地进行网格简化等操作; Simplygon[10]是一款商业网格简化工具，广泛应用于各种领域，其提供了多种简化算法，支持自动化简化流程，能够大幅度提高执行网格简化的效率; InstantUV是一款支持自动化 UV 映射的网格简化工具,可以大幅度提高 UV 映射和网格简化的效率，并能够快速生成高质量的网格。这些工具提供了丰富的网格处理和简化功能，并且支持多种网格数据格式的导入和导出。尽管存在诸多网格优化算法及网格处理工具，现有研究尚未考虑动画序列的优化，导致在复杂场景中渲染速度和存储开销的问题依然存在。

\subsection{贴图压缩研究现状}

贴图压缩针对 3D模型的贴图数据进行压缩。贴图文件需要转换为纹理才能被GPU 使用，因此贴图压缩可分为图片压缩及纹理压缩。Wu等提出JPEG编码照片集无损压缩方法，通过特征、空间和频域中的混合预测消除帧间和帧内几余，在不丢失任何信息的条件下进一步压缩了 JPEG 图片[11]。Mandeel 等介绍了目前通用图像压缩算法，在无损压缩模式下测试压缩比，实验结果表明对于单通道八位的图片JPEG XL 的平均压缩率最高[12]。图片压缩可减少图片尺寸，加快传输速度。纹理压缩对各类终端的纹理数据进行压缩，减少了显存使用量。Nah等提出一种基于亮度的 QuickETC2 技术，利用纹理块的亮度差异为每个块选择性的执行T或H压缩方法减少不必要的压缩开销[13]。贴图压缩能够降低贴图文件的冗余信息，为更好利用贴图压缩后的文件，可采用图像处理将单张贴图生成各种精度的图片，节省大量存储空间。Ye等提出一种基于双线性插值和小波变换的图像质量增强方法，利用直方图均衡和双线性插值增强图像，保留图像的细节质量，再利用小波重建对图像进行增强与双线性插值图像进行直方图匹配，得到更高精度的图像[14]。Zhang 等利用 mipmap生成多分辨率的纹理，基于网格可行性权重动态分配相应精度的纹理，保持高质量的渲染[15]。图片压缩和纹理压缩减少了数据传输时间，动态生成贴图方法减少了数据存储空间。现有研究未考虑结合两者，进一步加快贴图的获取速度。

\subsection{注视点渲染研究现状}

了解人类视觉系统的特性及其局限性有助于提高渲染质量和效率。眼睛与大脑共同构成了人类视觉系统，眼睛充当光学传感器，而大脑负责处理信息。光线进入眼睛后，首先穿过角膜折射，经过房水、虹膜、晶状体和玻璃体，最终到达视网膜，视网膜作为图像传感器接收视觉信息。

值得注意的是，并非所有进入眼睛的视觉信息都会传递至大脑。在视网膜层面，信息会经历大规模的压缩和预处理。视网膜上的感光器分布不均，其中视杆细胞和视锥细胞是两种主要的感光细胞类型。它们在形态、分布密度以及神经连接模式上存在显著差异[16]。视杆细胞在视网膜上分布较广，而视锥细胞则高度集中在视网膜的中央区域，即中央凹处，随着离中心的距离增大，视锥细胞的密度迅速降低[17]。这一分布特点导致视觉锐度在中央凹区域最高，并向外逐渐下降。注视点渲染技术便是基于这一生理特性，在中央凹区域提供高质量渲染，而在外围区域则降低渲染质量[18]。近年来，研究者在现有时空视觉模型的基础上，引入了多种感知优化，例如：Duinkharjav等人提出的颜色感知引导的显示功耗减少技术，通过优化颜色显示来降低VR和AR设备的功耗，同时保持视觉质量[19];Tariq等人引入的内容依赖性注视点技术，通过动态调整注视点渲染，优化了VR和AR中的视觉质量和计算效率[20];Tursun 和 Didyk引入的外周时空敏感性优化了头戴式显示器中外周视野的视觉体验，提高了渲染效率和视觉质量[21]; Krajancich引入的注意力机制通过优化注视点渲染，解决了头戴式显示器中因频繁眼动导致的视觉疲劳和计算资源浪费的问题[22]。然而现有研究并未充分考虑眼跳过程中的扫视抑制现象及眼跳后视觉锐度的恢复延迟，眼跳能够将人眼的注意焦点迅速转移至中央凹区域，从而弥补外围区域较低的视觉锐度，中央凹在视觉系统中具有最高的空间敏锐度和颜色感知能力[23]。在眼跳过程中，视觉输入会受到抑制，以隐藏视网膜的运动模糊，这种现象被称为扫视抑制[24]。需要注意的是，眼跳后视觉锐度的恢复并非瞬时完成，而是一个渐进的过程，通常伴随眼漂移和微扫视等不自觉的眼球运动[25]，这些运动会随着时间的推移逐渐提高对高空间频率的视觉敏感度[26]。
然而，现有的注视点渲染算法未能充分利用这一视觉恢复延迟特性，导致渲染效率未能得到最优化。

\subsection{视觉感知模型研究现状}

深入理解人眼在不同环境下的对比敏感度、空间频率响应以及注视点动态特性，对于设计高效的注视点渲染算法至关重要。近些年，研究人员基于过去的实验结果，提出了多种新颖的视觉感知模型来量化不同条件下人眼对细节的分辨能力。

在进行观察行为时，影响人眼视觉感知的因素有很多，单纯基于空间频率的渲染策略无法充分利用人眼的感知特性。 Ahumada 等人研究了视觉系统在空间和时间域中对信号的检测能力，提出了一个双通道空间-时间检测模型，该模型将视觉处理分为两个独立但相互关联的通道。环境光照强度和类型也会影响视觉敏感度，Tursun 等人提出了一种新的亮度对比度感知的注视点渲染技术，强调了亮度和对比度对人眼感知的影响，指出在不同光照条件下，用户的视觉敏感度会发生变化。Yi 等人也特别考虑了光照对人眼分辨能力的影响，提出了一种周围感知对比敏感度模型，专门用于描述在 HDR 显示环境中周围光照对对比敏感度的影响。
Wuerger 等人提出了一个空间-色彩对比敏感度模型，该模型综合了光照水平、空间频率和色彩对比度的影响，揭示了人眼在不同条件下的适应性和敏感度变化。Mantiuk 等人也提出了一种新的色彩对比敏感度模型，能够在极高亮度条件下有效预测人眼对色彩对比的感知。由于视网膜结构特性，目前大多数解决方案是将敏感度建模为偏心的函数。Watson 等人探讨了对比敏感度如何在不同视场和分辨率条件下变化，提出了一个多维感知模型，用于描述偏心度、分辨率和对比敏感度之间的关系。Krajancich 等人研究了在不同偏心度下，用户对空间和时间的闪烁融合能力，提出了一个偏心度依赖的感知模型，能够根据偏心度调整对闪烁的感知阈值。

尽管现有研究已经提出了多种新颖的视觉感知模型，但还未将综合了多维度因素的感知模型真正应用于注视点渲染技术当中。Kim 等人和 Peres 等人的研究仍是根据偏心度划分渲染区域，后续进行相应的渲染处理。Henriques 等人也使用偏心依赖的方式划分了三层渲染区域，但还考虑了光照强度对渲染结果的影响[30]。Surace 等人在注视点渲染中提出了两种用于简化网格的感知启发的可视性模型，分别关注空间
和时间频率的影响。Krajancich 等人提出了注意力感知的注视点渲染，以减少用于渲染的计算带宽。

综上所述，深入研究注视点渲染技术必然要关注视觉感知模型的建立，才能充分利用人类视觉特性，合理划分渲染区域或调整采样率等，从而设计更高质量且更高效的渲染方案。

\section{研究内容与主要工作}

基于以上分析的当前注视点渲染技术存在的挑战，为更有效地利用人类视觉系统特性，最大程度地减少渲染数据以及生成更高真实感的图像，以实现感知无损失的多分辨率渲染。本文基于现有的动画网格渲染方法，结合动态动画序列优化技术及眼跳后视觉锐度延迟恢复特性展开，提出了一种眼跳引导下的注视点渲染方法，该方法重点解决两个关键问题：一是通过网格数据精简与动画序列优化构建多维层次模型，以减少大规模动画网格数据带来的渲染延迟和存储开销；二是基于眼跳效应优化注视点渲染，充分考虑视觉锐度恢复的时间性变化，以降低渲染延迟。最终，研发出了眼跳引导下的虚拟现实动画网格高效渲染原型系统。具体研究内容包括以下几个方面：

（1）基于网格数据精简与动画序列优化的多维层次模型。针对动画网格数据规模过大显著增加渲染延迟和存储开销的问题，分析动画网格的组成，通过解耦静态网格与动画数据，采用网格合并和纹理压缩技术精简静态网格数据、优化关键帧和精简骨骼数据以优化动画序列，最终构建了一个适应不同分辨率渲染的多维层次模型，以降低渲染负担，为实现基于眼跳效应的注视点渲染技术奠定了基础。

（2）基于眼跳效应的注视点渲染方法。针对注视点渲染中忽略视觉锐度的时间性变化引发的延迟问题，本文研究了扫视抑制现象及眼跳后视觉锐度延迟恢复的特性。基于此，提出了一种基于眼跳效应的注视点渲染方法，通过在扫视过程中及眼跳后的短时间内降低渲染分辨率，从而减少潜在计算和数据消耗，提高渲染效率。

（3）设计并实现眼跳引导下的虚拟现实中动画网格的高效渲染原型系统。结合上述两个技术方法，本文设计并实现了眼跳引导下的虚拟现实中动画网格的高效渲染原型系统，有效解决了虚拟现实中动画网格渲染延迟与存储开销的问题。

\section{论文组织结构安排}
xxxxx